\documentclass[11pt, a4paper]{article}

% ── Packages ──────────────────────────────────────────────────────────
\usepackage[utf8]{inputenc}
\usepackage[T1]{fontenc}
\usepackage{lmodern}
\usepackage[margin=1in]{geometry}
\usepackage{amsmath, amssymb, amsthm}
\usepackage{booktabs}
\usepackage{longtable}
\usepackage{graphicx}
\usepackage{hyperref}
\usepackage{xcolor}
\usepackage{listings}
\usepackage{tcolorbox}
\usepackage{polya}

\hypersetup{
  colorlinks=true,
  linkcolor=polyablue,
  urlcolor=polyablue,
  citecolor=polyablue,
}

% ── Code listing style ────────────────────────────────────────────────
\lstdefinestyle{polya-python}{
  language=Python,
  basicstyle=\ttfamily\small,
  keywordstyle=\color{polyablue}\bfseries,
  commentstyle=\color{polyagray}\itshape,
  stringstyle=\color{polyagreen},
  numbers=left,
  numberstyle=\tiny\color{polyagray},
  numbersep=8pt,
  frame=single,
  framerule=0.4pt,
  rulecolor=\color{polyagray},
  breaklines=true,
  showstringspaces=false,
  tabsize=4,
}
\lstset{style=polya-python}
\title{Causal Inference -- Job Training Program}
\author{uber-polya}
\date{2026-02-26}

\begin{document}

\maketitle
\tableofcontents
\newpage

% ── Phase A: Problem Formulation ─────────────────────────────────────
\section{Problem Formulation}

\begin{polyamodel}

\textbf{Problem Type}: Find \hfill
\textbf{Domain}: Causal Inference


\end{polyamodel}

% ── Universe ──────────────────────────────────────────────────────────
\subsection{Universe}

\begin{itemize}
  \item $Naive: \{estimate, bias\}$
  \item $Propensity Score Matching: \{att, propensity_treated_mean, propensity_control_mean, error\}$
  \item $Difference In Differences: \{estimate, error\}$
  \item $Doubly Robust: \{ate, error\}$
\end{itemize}

% ── Variables ─────────────────────────────────────────────────────────
\subsection{Variables}

\begin{longtable}{lll}
\toprule
\textbf{Symbol} & \textbf{Meaning} & \textbf{Type / Range} \\
\midrule
\endhead
$x$ & decision variable & $\mathbb{{R}}$ \\
\bottomrule
\end{longtable}

% ── Structure ─────────────────────────────────────────────────────────
\subsection{Structure}

- **Treated**: 500 employees who participated in the training program (approximately)
- **Control**: 500 employees who did not participate (approximately)
- **Covariates**: education\_years, experience\_years, age (all affect both participation and salary)
- **Confounding**: Higher-education, more-experienced employees self-select into training AND earn more
- **Question**: What is the causal effect of the training program on salary?
- **Ground truth**: The data is synthetic with a true treatment effect of \$5,000

% ── Mapping ───────────────────────────────────────────────────────────
\subsection{Real-World Mapping}

\begin{longtable}{ll}
\toprule
\textbf{Real-World Concept} & \textbf{Mathematical Object} \\
\midrule
\endhead
problem input & $instance parameters$ \\
\bottomrule
\end{longtable}

% ── Constraints ───────────────────────────────────────────────────────
\subsection{Constraints}

\begin{enumerate}
  \item $Selection bias**: Treated and control groups differ systematically on covariates, inflating naive estimates$
  \hfill \textit{()}
  \item $Propensity score**: P(treatment | covariates)$
  \hfill \textit{(summarizes confounding into a single number for matching)}
  \item $ATE vs ATT**: Average Treatment Effect (population) vs Average Treatment Effect on the Treated (those who participated)$
  \hfill \textit{()}
  \item $Confounding**: When a variable affects both treatment assignment and outcome$
  \hfill \textit{()}
  \item $Doubly robust estimation**: Combines propensity model with outcome model; consistent if either model is correct$
  \hfill \textit{()}
\end{enumerate}

% ── Objective / Claim ─────────────────────────────────────────────────
\subsection{Claim}


% ── Approach ──────────────────────────────────────────────────────────
\subsection{Approach}

\begin{description}
  \item[Suggested approach:] Naive + Propensity Score Matching + Difference-in-Differences + Doubly Robust (AIPW)
  \item[Complexity class:] 
  \item[Available tools:] Naive + Propensity Score Matching + Difference-in-Differences + Doubly Robust
\end{description}
% ── Phase B: Solution ────────────────────────────────────────────────
\section{Solution}

\begin{polyaresult}

\textbf{Answer}: Solution computed


\textbf{Optimal}: No / Unknown \hfill
\textbf{Feasible}: No
\textbf{Algorithm}: Naive + Propensity Score Matching + Difference-in-Differences + Doubly Robust (AIPW) \quad $$

\textbf{Time}: 0.009365651989355683s


\end{polyaresult}

% ── Solution Details ──────────────────────────────────────────────────
\subsection{Solution Details}

Solution computed successfully.

% ── Verification ──────────────────────────────────────────────────────
\subsection{Verification}

\begin{longtable}{lcl}
\toprule
\textbf{Check} & \textbf{Status} & \textbf{Detail} \\
\midrule
\endhead
Naive Biased & \passmark & Passed \\
Psm Close & \passmark & Passed \\
Did Close & \passmark & Passed \\
Dr Close & \passmark & Passed \\
Psm Closer Than Naive & \passmark & Passed \\
Dr Closer Than Naive & \passmark & Passed \\
Propensity Overlap & \passmark & Passed \\
Treatment Balance & \passmark & Passed \\
All Passed & \passmark & Passed \\
\bottomrule
\end{longtable}

% ── Phase C: Interpretation ──────────────────────────────────────────
\section{Interpretation}

% ── The Question & The Answer ─────────────────────────────────────────
\subsection{The Question}
- **Treated**: 500 employees who participated in the training program (approximately)
- **Control**: 500 employees who did not participate (approximately)
- **Covariates**: education\_years, experience\_years, age (all affect both participation and salary)
- **Confounding**: Higher-education, more-experienced employees self-select into training AND earn more
- **Question**: What is the causal effect of the training program on salary?
- **Ground truth**: The data is synthetic with a true treatment effect of \$5,000.

\subsection{The Answer}

\begin{polyainsight}[title={Bottom Line}]
A feasible solution was found.
\end{polyainsight}

% ── What This Means ───────────────────────────────────────────────────
\subsection{What This Means}

- Naive difference: Overestimates the effect (well above \$5,000) due to selection bias
- Propensity Score Matching (ATT): Recovers approximately \$5,000
- Difference-in-Differences: Recovers approximately \$5,000
- Doubly Robust (AIPW): Recovers approximately \$5,000
- Verification: All 8 checks pass -- naive is biased, causal methods are closer to truth

% ── Sensitivity Analysis ──────────────────────────────────────────────

% ── Figures ───────────────────────────────────────────────────────────

% ── Recommendations ───────────────────────────────────────────────────
\subsection{Recommendations}

\begin{enumerate}
  \item Review the model assumptions to ensure real-world fidelity.
\end{enumerate}

% ── Limitations ───────────────────────────────────────────────────────
\subsection{Limitations}

\begin{polyawarnbox}
\begin{itemize}
  \item Model assumes deterministic parameters; real-world variability not captured.
  \item Solution quality depends on the accuracy of input data.
\end{itemize}
\end{polyawarnbox}

% ── Appendix ─────────────────────────────────────────────────────────

\end{document}